\subsection*{92. Candidate B: Engineering Colleges Seeking Applied Capabilities}

Imagine a private engineering college with students who have learned conventional coursework but struggle to design and defend meaningful applied projects. Administrators want an industry-relevant elective, project lab, faculty support, or a structured applied experience. The college may have reputational and student-outcome incentives, but also constrained budgets, faculty workload, timetables, curriculum approval processes, and procurement requirements.

Potential Bayes capability: curriculum and technical teaching, ability to create rigorous assessments and applied projects, and some qualitative exposure to college selling. Potential wedge: a tightly scoped, faculty-enabled applied AI elective or lab with measurable student work and clear service boundaries.

Positive conversations with a principal or training head do not prove that the college will sign, pay promptly, assign time, supply computers, or support quality standards. The real beachhead might be 1 category of institution with a particular academic and purchasing structure, not ``colleges across India.''

\subsection*{93. Narrow Candidate B by Purchase Mechanism}

A college promising a guest lecture, a department willing to allocate an elective, an autonomous institution able to approve curricular changes, and a placement office needing a seasonal project workshop are not equivalent buyers. They have different budgets, contracts, decision rights, student prerequisites, and evidence requirements.

A useful candidate might be: ``private engineering institutions with flexible electives, a named internal faculty coordinator, an approved budget for applied AI, and an upcoming semester window.'' Whether this group can be reached and sold within the founder's available hours must be tested.

The initial offer might be a bounded applied module with a published delivery calendar, clearly allocated faculty roles, project assessments, no unearned affiliation claims, and milestone payments. Avoid treating a 1-day placement workshop as proof of demand for a semester-long curricular partnership.

\subsection*{94. Strategic Assets and Hidden Liabilities of the College Route}

Potentially valuable accumulated assets include instructional rubrics tested across ability levels, faculty-enablement procedures, relationships with academic departments, applied laboratory templates, and employer-facing student artifacts. These can align with the institution's long-term vision.

Risks include travel and calendar conflicts with the founder's job, price pressure, slow collections, requirements to train very large heterogeneous groups, negotiated customizations, and administrators pushing for impressive completion metrics without rigorous standards. A big contract can generate revenue while eroding the exact academic reputation Bayes hopes to build.

A good first college beachhead would involve a manageable academic calendar and 1 sponsor capable of operating alongside Bayes, not dependence on the founder doing every class personally.

\subsection*{95. Candidate C: An Enterprise Team Facing a Specific AI Gap}

An enterprise AI capability buyer may not want a ``course'' at all. They may need a team to produce, evaluate, or maintain 1 family of AI-enabled workflows. Their benchmark may be delivery time, error reduction, reliability, cost control, or internal review quality.

Potential wedge: a paid diagnostic combining skill assessment with a non-confidential workflow or architecture exercise, followed by a limited training-and-implementation sprint. The natural proof may be documented pre/post performance on relevant tasks and the client's assessment of operational impact, not generic learner satisfaction.

The economic buyer and training participants differ. An enthusiastic engineering manager may lack L\&D budget; an L\&D manager may fund workshops but lack ownership of the technical problem. The beachhead must contain an identifiable sponsor who can act.

\subsection*{96. Narrow Candidate C to a Purchasable Business Situation}

``Enterprise AI upskilling'' is too broad. Possible narrower candidates include engineering teams about to deploy customer-facing AI features, analytics units adopting natural-language data interfaces, or product teams with unreliable LLM evaluation practices.

For each, ask whether Bayes has domain expertise, whether real tasks can be discussed within confidentiality rules, whether a small engagement has an economic buyer, and whether the promised impact can be measured without invented ROI. B2B technical training can become deeply bespoke. This may be profitable but founder-dependent unless the methodology is productized and the custom boundary is controlled.

\subsection*{97. A Transparent Comparison of the 3 Initial Routes}

The following is a \textbf{qualitative hypothesis matrix}, not an evidence-backed ranking or recommendation to commit today:

\begin{table*}[t]
\centering
\small
\begin{tabularx}{\textwidth}{l p{4.6cm} p{4.6cm} X}
\toprule
\textbf{Dimension} & \textbf{D2C Professional Studio} & \textbf{College Applied Module} & \textbf{Enterprise Team Diagnostic} \\
\midrule
End user & Working engineer & Students and department faculty & Enterprise engineering employees \\
Economic buyer & Usually individual learner & College or department management & Engineering manager or L\&D sponsor \\
Core job & Independent technical capability & Curriculum-integrated applied lab & Team capability tied to work deliverable \\
Initial proof & Public sample task and rubric & Pilot academic delivery and artifacts & Diagnostic audit and task improvement \\
Access hypothesis & Practitioner circles and peer referrals & Direct outreach and academic networks & Targeted executive introductions \\
Payment horizon & Short, but highly price-sensitive & Tied to academic semester budgets & Procurement and budget cycles vary \\
Cost risk & Individual reviews and acquisition & Coordination, travel, delayed cash & Customization and long sales cycles \\
Institutional asset & Assessment rubrics, learner proof & Curricular modules, faculty processes & Industry cases, production workflows \\
Transfer limit & Learner proof does not persuade deans & College proof does not prove seniority & Enterprise ROI does not prove pedagogy \\
Most vital unknown & Will engineers pay and sustain effort? & Will deans fund a rigorous pilot? & Will sponsors pay for measured gains? \\
\bottomrule
\end{tabularx}
\caption{Comparative qualitative hypothesis matrix of Bayes Institute's 3 candidate entry routes.}
\label{tab:ch7-bayes-routes-matrix}
\end{table*}

All 3 could be valuable eventually. The immediate task is to determine which has the strongest evidence-adjusted near-term entry under a 20-hour founder schedule and a very small budget.

\subsection*{98. Provisional Prioritization for Bayes Institute}

Given the founder's technical background, ability to teach personally, desire to build direct academic identity, avoidance of heavy upfront fixed costs, and need for repeatable cash flow, I would make a \textbf{narrow paid technical learning or assessment offer to working engineers the default first test}, not a large immediate 6-month launch. This is a recommendation for experiment sequence, not a conclusion that D2C is the winning business.

In parallel, I would allow a strictly bounded college or enterprise opportunity to be tested if a real, funded buyer with a clear scope emerges through existing relationships. I would not create 3 full marketing funnels, brands, and curricula at once.

Why the provisional D2C preference? The founder can directly demonstrate core competence, observe learning, control quality, publish relevant educational evidence with permission, and test payment without waiting for a university procurement committee. The risks are real (self-pay demand, competition, schedules, and labor-intensive feedback) and must be tested quickly. If evidence is weak but a college buyer presents a genuine funded pilot with favorable economics, the ranking can reverse.
